% Compile with pdflatex, xelatex, or tectonic:
%   tectonic report.tex      (recommended; self-contained)
%   pdflatex report.tex      (run twice to resolve references)
\documentclass[11pt]{article}

\usepackage[T1]{fontenc}
\usepackage{lmodern}
\usepackage{graphicx}
\usepackage{booktabs}
\usepackage{geometry}
\geometry{a4paper, margin=2.5cm}
\usepackage{hyperref}
\hypersetup{colorlinks=true, linkcolor=blue, urlcolor=blue}

\title{ERA Official Demo: Stage-1 Reproduction Report}
\author{Reproduction Experiment Log}
\date{\today}

\begin{document}
\maketitle

\section{Overview}
This report documents a reproduction of the \textbf{official minimal demo} shipped
with \href{https://github.com/google-research/era}{\texttt{google-research/era}},
the open-source code for the Nature paper \emph{An AI system to help scientists
write expert-level empirical software}.

An important clarification: this is \textbf{not} the full scRNA-seq or COVID
scientific experiment from the paper. It is the official minimal, runnable demo
included in the repository. Its purpose is to first get the entire ERA pipeline
working locally end-to-end (LLM sampling, sandboxed execution, tree search,
scoring, and iterative improvement), confirming that the system runs correctly
and exhibits the expected ``gradual self-improvement'' behaviour, as a Stage-1
foundation for the larger experiments to follow.

\section{The Task}
The demo is a Kaggle / California Housing--style \textbf{regression task}: given
housing-related features, predict a continuous target value. ERA first evaluates
an initial baseline solution and then runs a tree search that repeatedly
generates, executes, and improves candidate programs.

\section{The Score}
The scoring metric is the \textbf{negative root-mean-squared error
(negative RMSE)}, i.e. $-\mathrm{RMSE}$. Because of the negative sign, a
\textbf{higher score (closer to $0$) means a better model}.

\section{Key Results}
The core results of this reproduction are summarised below.

\begin{table}[h]
\centering
\begin{tabular}{lc}
\toprule
\textbf{Stage} & \textbf{Best score (negative RMSE, higher is better)} \\
\midrule
Initial baseline        & $-0.7339448894714724$ \\
10-iteration best       & $\approx -0.5785$ \\
30-iteration best       & $-0.5775746202235905$ \\
\bottomrule
\end{tabular}
\caption{Best score of the ERA official demo at each stage.}
\end{table}

\section{Search Progress}
Figure~\ref{fig:progress} shows the best-so-far score as a function of the search
step during ERA tree search. Step~0 is the initial linear-regression baseline; a
large jump occurs immediately at step~1, after which the search makes small,
steady improvements on gradient-boosting / feature-engineering solutions before
quickly reaching a plateau.

\begin{figure}[h]
\centering
\includegraphics[width=0.85\textwidth]{progress}
\caption{ERA demo search progress: best-so-far score (negative RMSE) versus search
step. The red dashed line is the initial baseline, the green dotted line is the
10-iteration best, and the purple dash-dotted line is the 30-iteration best.}
\label{fig:progress}
\end{figure}

\section{Interpretation}
\begin{itemize}
  \item \textbf{The ERA main loop was successfully reproduced.} The system
        completes the full closed loop of
        sample $\rightarrow$ sandboxed execution $\rightarrow$ scoring
        $\rightarrow$ selection and iteration, and exhibits the expected
        self-improvement behaviour.
  \item \textbf{Solution quality improved substantially.} The search started from
        a simple linear-regression baseline (score $\approx -0.7339$) and
        automatically evolved into a solution centred on gradient boosting
        (\texttt{HistGradientBoostingRegressor}), with extensive feature
        engineering (including location/distance features) and an ensemble of
        three gradient-boosting models (score $\approx -0.5776$).
  \item \textbf{It plateaus quickly.} Thirty iterations improved only marginally
        over ten iterations (roughly $-0.5785 \rightarrow -0.5776$), indicating
        that the search space of this toy task is small and that ERA approaches
        the plateau quickly, with rapidly diminishing returns from further
        iterations.
\end{itemize}

\section{Next Step}
The next experiment is to compare, on the same task,
\textbf{ERA tree search} against \textbf{best-of-$N$ independent LLM sampling}.
Under an equal LLM-call budget, we will evaluate the final score and sample
efficiency of both approaches, in order to measure the actual gain of ERA's
``tree search + iterative improvement'' over plain independent parallel sampling.

\end{document}
